\begin{enumerate}
    \item[\textbf{Problema 4.1.}] Mostre que a eq. (4.7) pode ser obtida substituindo $ix$ por $x$ na eq. (4.6).
    \item[\textbf{Problema 4.2.}] Determine os primeiros quatro termos das expansões de Taylor de $\tan x$ e $\cot x$ sobre $x = 0$.
    \item[\textbf{Problema 4.3.}] Determine os quatro primeiros termos das expansões de Taylor de $\tanh x$ e $\coth x$ sobre $x = 0$.
    \item[\textbf{Problema 4.4.}] Use a análise dimensional para determinar como o parâmetro catenário, $c$, está relacionado ao componente horizontal constante da tensão do cabo, $T_0$, e seu peso por unidade de comprimento (ou peso unitário), $\gamma$.
    \item[\textbf{Problema 4.5.}] Quanta curvatura de fita (sag) é permitida para medir uma distância de 50 m com precisão de 5%? Dentro de 2%?.
    \item[\textbf{Problema 4.6.}] O que pesa um corpo que pesa 10 N na superfície da Terra a uma altura de 10 m? No pico do Monte Everest? (Dica: você pode ter que procurar alguns fatos sobre o nosso planeta!).
    \item[\textbf{Problema 4.7.}] De acordo com a eq. (4.30), em que altitude o peso de 10 N na superfície da Terra cairia para 9 N? Para 5 N?.
    \item[\textbf{Problema 4.8.}] Compare os resultados obtidos no Problema 4.7 com os resultados mais exatos obtidos usando a eq. (4.29).
    \item[\textbf{Problema 4.9.}] O que pesa um corpo que pesa 10 N na superfície da Terra na superfície da lua? Na superfície do planeta Plutão? Na superfície do planeta Marte? (Dica: você pode ter que procurar alguns fatos sobre o ambiente de nossos planetas!).
    \item[\textbf{Problema 4.10.}] Se o potencial gravitacional correspondente à lei da gravitação de Newton (eq. (4.25)) é dado por $V_g = -GM_Em/R$, encontre a expressão exata que define esse potencial em função da altitude, $z$, a partir da superfície da Terra.
    \item[\textbf{Problema 4.11.}] Escreva uma expansão binomial dos resultados do Problema 4.10 para determinar a energia potencial acima da superfície da Terra para a primeira ordem em $z$.
    \item[\textbf{Problema 4.12.}] Preencha os elementos ausentes da tabela a seguir para a ordem de dois termos: $\sin x$, $\cos x$, $1 - \sin x$, $1 - \cos x$.
    \item[\textbf{Problema 4.13.}] Preencha os elementos ausentes da tabela a seguir para a ordem de dois termos: $\sinh x$, $\cosh x$, $1 - \sinh x$, $1 - \cosh x$.
    \item[\textbf{Problema 4.14.}] Desenvolva um coeficiente de expansão de volume, $\beta$, para um sólido de comprimento $L_0$, largura $W_0$ e altura $H_0$, que seja paralelo ao coeficiente de superfície, $\gamma$, da eq. (4.37).
    \item[\textbf{Problema 4.15.}] A que diferença de temperatura um sólido de alumínio teria de ser submetido para o coeficiente de superfície de expansão produzir erros de 1% na mudança de área em comparação com a mudança exata de área?.
    \item[\textbf{Problema 4.16.}] A que diferença de temperatura um sólido de alumínio teria que ser submetido ao coeficiente de expansão de volume para produzir erros de 1% na mudança de área em comparação com a mudança exata de área?.
    \item[\textbf{Problema 4.17.}] Arredonde cada um dos seguintes números para dois (2) algarismos significativos: (a) 5,237 (b) 0,82549 (c) 81,356 (d) $\pi$ (e) 6,2305 (f) 0,0428 (g) 10,45 (h) 4,035.
    \item[\textbf{Problema 4.18.}] Arredonde cada um dos seguintes números para três (3) algarismos significativos: (a) 5,237 (b) 0,82549 (c) 81,356 (d) $\pi$ (e) 6,2305 (f) 0,0428 (g) 10,45 (h) 4,035.
    \item[\textbf{Problema 4.19.}] Complete as seguintes multiplicações e expresse os resultados para o número correto de algarismos significativos: (a) $(6,28 \times 10^3) \times 2,712$ (b) $43,32 \times 0,3$ (c) $928 \times 4,23$.
    \item[\textbf{Problema 4.20.}] 99,9 e 100,1 têm o mesmo número de algarismos significativos? Explique sua resposta.
    \item[\textbf{Problema 4.21.}] Estime os intervalos dentro dos quais cada um dos seguintes números se encontra: (a) 7,7 (b) 7,70 (c) 1200 (d) $1,200 \times 10^{-3}$.
    \item[\textbf{Problema 4.22.}] Em que porcentagem o período de um pêndulo mudaria se seu comprimento fosse reduzido pela metade? Se fosse reduzido em um terço? Se o comprimento fosse reduzido para um terço de seu comprimento original?.
    \item[\textbf{Problema 4.23.}] Explique por que o período do pêndulo aumenta em 145% na lua.
    \item[\textbf{Problema 4.24.}] Como o período de um pêndulo mudaria, em comparação com seu valor na Terra, se o pêndulo estivesse em Marte? Em Plutão?.
    \item[\textbf{Problema 4.25.}] Como o período de um pêndulo mudaria em função de sua altura, $h$, acima da superfície da Terra? (Dica: A variação da aceleração gravitacional $g$ pode ser representada como uma função de $h$ a partir da lei da atração gravitacional de Newton).
    \item[\textbf{Problema 4.26.}] Desenhe dois alvos circulares de arco e flecha e use-os para representar os padrões de acerto de (a) um arqueiro que é exato (accurate), mas não preciso (precise); e (b) um arqueiro que é preciso, mas não exato.
    \item[\textbf{Problema 4.27.}] Verifique as formas finais das eqs. (4.51) e (4.52).
    \item[\textbf{Problema 4.28.}] Verifique as equações para $m$ e $b$ dadas nas eqs. (4.53) e (4.54).
    \item[\textbf{Problema 4.29.}] Discuta e explique as diferenças dimensionais entre as eqs. (4.53) e (4.54).
    \item[\textbf{Problema 4.30.}] Verifique os termos na terceira e quarta colunas da Tabela 4.3, bem como as somas de todas as quatro colunas.
    \item[\textbf{Problema 4.31.}] Verifique os cálculos de $m$ e $b$ encontrados nos resultados da Tabela 4.3.
    \item[\textbf{Problema 4.32.}] Determinar o desvio-padrão dos dados apresentados no quadro 4.4.
    \item[\textbf{Problema 4.33.}] Desenhe um histograma para os dados na Tabela 4.4 com 10 intervalos de largura de 4 dB.
    \item[\textbf{Problema 4.34.}] Mostre que o quadrado da variância amostral da eq. (4.57) pode ser lançado na forma alternativa $s^2 = \frac{1}{(n-1)} \left[ \sum_{i=1}^n x_i^2 - \frac{1}{n} \left(\sum_{i=1}^n x_i\right)^2 \right]$.
    \item[\textbf{Problema 4.35.}] Estime o erro cometido na aproximação de $y(x) = \sin x$ com uma fórmula de Taylor para $n = 4$ calculando o restante $R_5$.
    \item[\textbf{Problema 4.36.}] As afirmações de que $\sin x \simeq x$ e $\tan x \simeq x$ produzem aproximações semelhantes? Confirme e explique sua resposta.
    \item[\textbf{Problema 4.37.}] As leituras de um voltímetro analógico antiquado (ele tem mostradores, não leituras digitais!) estão sujeitas a algum erro sistemático quando todas as suas leituras são muito grandes. Descobriu-se que a magnitude do erro varia linearmente de 1 V com uma leitura de discagem de 5 V a 4 V com uma leitura de discagem de 80 V. (a) Quais são as tensões corretas para leituras de discagem de 80, 100, 50, 1, 35 e 10 V? (b) Qual é o erro percentual para cada uma das seis (6) leituras da parte (a)?.
    \item[\textbf{Problema 4.38.}] (a) É possível dispor de um conjunto de medições precisas, mas não exatas? Explicar. (b) É possível ter um conjunto de medições exatas, mas não precisas? Explicar.
    \item[\textbf{Problema 4.39.}] (a) Escreva a expansão da série de Taylor para $e^x$ sobre $x = 0$. (b) Calcule $e^{0,5}$ para cinco algarismos significativos usando os quatro primeiros termos da série encontrada na parte (a).
    \item[\textbf{Problema 4.40.}] (a) Que erro percentual foi incorrido no cálculo da parte (b) do Problema 4.39 se o valor real de $e^{0,5}$ for 1,6487? (b) Use o resto de Taylor (eq. (4.5)) para calcular o erro em $e^{0,5}$ após apenas quatro termos. O erro calculado na parte (a) deste problema é aceitável? Explicar.
    \item[\textbf{Problema 4.41.}] Avalie a seguinte função manualmente (sem calculadoras ou computadores, por favor) para $x = 4$: $(1 + 2/x)^{1/4}$.
    \item[\textbf{Problema 4.42.}] Como é que um observador sabe quando é suficiente, que foram feitas medições suficientes?.
    \item[\textbf{Problema 4.43.}] Faça uma lista de cinco novos exemplos (ou seja, não encontrados no texto) de erros sistemáticos.
    \item[\textbf{Problema 4.44.}] Faça uma lista de cinco novos exemplos (ou seja, não encontrados no texto) de erros aleatórios.
    \item[\textbf{Problema 4.45.}] A resistência de um resistor, $R$, é feita passando várias correntes, $I$, através dele e medindo as quedas de tensão correspondentes, $V$, e correntes com medidores analógicos imprecisos. Os dados resultantes são: $x_i = V$ (V): 10, 20, 30, 40, 50, 60, 70, 80; $y_i = I$ (A): 0.8, 1.1, 2.5, 4.2, 4.3, 4.7, 5.8, 6.4. (a) Que tipos de erro serão encontrados nos dados? (b) Supondo que $V = IR$, plote os dados (manualmente!) e estime visualmente (globo ocular) a linha de melhor ajuste para esses dados.
    \item[\textbf{Problema 4.46.}] Use o método dos mínimos quadrados para traçar uma curva $V$ versus $I$ para os dados do Problema 4.45. Como ele se compara com o resultado visual (globo ocular) do Problema 4.45?.
    \item[\textbf{Problema 4.47.}] Os dados apresentados a seguir compreendem 100 leituras dos níveis de ruído efetuados a 6 milhas de distância de um aeroporto, efetuadas tarde da noite a intervalos de 15 segundos. Encontre a média, a mediana e o desvio padrão desses dados.
    \item[\textbf{Problema 4.48.}] Os números estrelados nos dados do Problema 4.47 são leituras feitas enquanto uma aeronave estava voando diretamente acima. Se esses dados forem excluídos, quais são a média, a mediana e o desvio padrão dos 88 pontos de dados restantes?.
    \item[\textbf{Problema 4.49.}] Traçar (a) um histograma de todos os dados do problema 4.47 e (b) uma curva contínua do número de leituras em função do nível de ruído medido.
    \item[\textbf{Problema 4.50.}] Determinar a aproximação de campo distante da função $f(r)$ dada abaixo como uma expansão binomial para valores de $r \gg a$: $f(r) = \sqrt{a^2 + r^2}$.
    \item[\textbf{Problema 4.51.}] O potencial elétrico, $V_e$, à distância, $r$, ao longo do eixo de revolução de um disco de raio $a$ é dado por $V_e = \frac{q}{2\pi a^2 \varepsilon_0}(\sqrt{a^2 + r^2} - r)$, onde $q$ é a carga total que é distribuída uniformemente sobre a superfície do disco e $\varepsilon_0$ é a constante de permissividade. Usando os resultados do Problema 4.50, encontre uma aproximação de campo distante para o potencial elétrico para valores de $r \gg a$.
    \item[\textbf{Problema 4.52.}] Compare o número mínimo de termos mantidos nas expansões binomiais das soluções para os problemas 4.50 e 4.51. Esses números são iguais ou não? Por que esses números são iguais ou não?.
    \item[\textbf{Problema 4.53.}] Suponha que precisamos calcular a extensão radial ou deflexão $w$ de um balão esférico muito fino, o que significa que o raio da esfera se estende de $R$ a $R + w$ à medida que o balão é pressurizado. É feito de material elástico. Um colega encontra um livro didático que mostra uma fórmula para a pressão, $p$, que parece razoável: $w/R = pR/Eh$, onde $h$ é a espessura da parede do balão e $E$ é o módulo do material do qual a esfera é feita. Esta equação é dimensionalmente consistente?.
    \item[\textbf{Problema 4.54.}] Analise o comportamento limite da equação apresentada no Problema 4.53 à medida que a pressão, o módulo, o raio e a espessura vão para zero e se tornam infinitamente grandes. Esse comportamento limite está de acordo com sua estimativa intuitiva do que deve acontecer?.
    \item[\textbf{Problema 4.55.}] Use a equação do Problema 4.53 para derivar uma estimativa da magnitude da pressão, $p$, como uma fração do módulo, $E$. Estimar a fração de pressão para uma esfera de paredes finas, para a qual $h/R \ll 1$.
\end{enumerate}
